

\begin{document}
	
	% RESET all counters at the beginning
	\setcounter{section}{0}
	\setcounter{subsection}{0}
	\setcounter{subsubsection}{0}
	\setcounter{paragraph}{0}
	
	% Depth for numbering and TOC
	\setcounter{secnumdepth}{1}  % Number sections only
	% Part in TOC: footnotesize, bold, NO page break
	\makeatletter
	\renewcommand*\l@part[2]{%
		\ifnum \c@tocdepth >-2\relax
		\addpenalty{-\@highpenalty}%
		\addvspace{0.8em \@plus\p@}%
		{\leftskip 0em \relax
			\rightskip \@tocrmarg
			\parfillskip -\rightskip
			\parindent 0em \relax\@afterindenttrue
			\interlinepenalty\@M
			\leavevmode
			{\footnotesize\bfseries #1}\nobreak
			\leaders\hbox{$\m@th\mkern \@dotsep mu\hbox{}\mkern \@dotsep mu$}\hfill
			\nobreak\hb@xt@\@pnumwidth{\hss #2}\par}%
		\addvspace{0.2em \@plus\p@}%
		\nobreak
		\fi}
	\makeatother
	
	% Chapter in TOC: footnotesize, bold
	\renewcommand{\cftchapfont}{\footnotesize\bfseries}
	\renewcommand{\cftchappagefont}{\footnotesize\bfseries}
	\setlength{\cftbeforechapskip}{0.3em}
	
	% Only Chapters in TOC (no Sections/Subsections)
	\setcounter{tocdepth}{0}
	

	\begin{titlepage}
		\centering
		\vspace*{2cm}
	{\Huge\textbf{Fundamentale Fraktal-Geometrische Feldtheorie (FFGFT) oder T0-Theorie: Zeit-Masse-Dualität}}\\[1cm]
	{\Large Teil 6: Holografie, HLV-Brücke und die Struktur der Zeit}\\[2cm]
		{\large Johann Pascher}\\[1cm]
		{\large 2026}
		\vfill
	\end{titlepage}

	\frontmatter
	\pagestyle{fancy}
	\makeatletter
	\let\ps@plain\ps@fancy
	\let\ps@empty\ps@fancy
	\makeatother

	\mainmatter
	\pagestyle{fancy}
	\tableofcontents

	\chapter{Einleitung zu Band 6}
	% =============================================================================
% EINLEITUNG ZU BAND 6: HOLOGRAFIE, HLV-BRÜCKE UND DIE STRUKTUR DER ZEIT
% (Bd. 6 von 8; Fortsetzung von Bd. 5)
% =============================================================================

\section*{Sechster Band der Dokumentensammlung}

Mit diesem Band wird die Sammlung der Einzeldokumente zur T0-Theorie / FFGFT
über die ursprünglich fünf Bände hinaus fortgeführt. Seit dem Abschluss von
Band~5 (Dok.~262, Mai 2026) ist der Korpus um mehr als 120 Dokumente
gewachsen. Sie werden in den Bänden 6, 7 und 8 in numerischer Reihenfolge
aufgenommen; die Grenze jedes Bandes ist so gelegt, dass ein inhaltlicher
Schnitt mit dem Umfangsmaß der Druckausgaben zusammenfällt.

\subsection*{Inhalt und Schwerpunkt}

Band~6 umfasst die Dokumente 263 bis 310 aus den Monaten Juni und Juli 2026.
Den roten Faden bildet die Frage, wie die kompakte Geometrie $T^4$ in die
beobachtete dreidimensionale, zeitlich offene Welt übergeht und was dabei an
Information, Phase und Zeitstruktur erhalten bleibt. Am Anfang stehen das
holografische Prinzip in fraktaler Lesart, die kosmologische Entartung
zwischen $\Lambda$CDM und FFGFT sowie die Herleitung der CMB-Peak-Verhältnisse
aus dem $T^4$. Es folgt die Projektionskette vom $T^4$ zum Beobachtungsraum,
die über eine Reihe von Vergleichs- und Brückendokumenten zum
Helix--Light--Vortex-Ansatz (HLV) bis zur gerechneten 3-gegen-6-Brücke in
Dok.~285 führt.

Der dritte Abschnitt behandelt den dynamischen Sektor: kinetische Energie als
$\xi$-fraktalen Anteil des statischen Verhältnisses $\tilde T\cdot m=1$, die
Trennung von Betrag und Phase in den vier Fermion-Sektoren und die doppelte
Herkunft des Winkels $\theta=2/9$, dynamisch (Dok.~291) und ikosaedrisch
(Dok.~293). Die Dokumente 295 bis 307 widmen sich der Zeit und der
Information: die Zeit-Windung als logarithmische Spirale, der Umkehrtest, der
Information als Ausgang und nicht als Grundeinheit ausweist, die native
Zeit-Energie-Reziprozität und die Verortung der Zeit im Zustandsraum. Den
Abschluss bilden der kosmische Sektor und das Skalenanker-Problem im Vergleich
mit $\Lambda$CDM sowie die Resonanzgeometrie als übergeordnete Referenz.

\subsection*{Gliederung}

Die 44 Dokumente sind in fünf Teile gegliedert: Holografie, kosmologische
Entartung und Formalismus-Korrespondenzen (Dok.~263--269);
Dimensionswechsel, HLV-Brücke und Skalenrekursion (Dok.~270--285);
dynamischer Sektor, Betrag und Phase, $\theta=2/9$ (Dok.~286--294); Zeit und
Information (Dok.~295--307); kosmischer Sektor und Resonanzgeometrie
(Dok.~308--310).

\subsection*{Zum Lesen der Dokumente}

Wie in den früheren Bänden ist jedes Dokument eigenständig und wird im Text
mit seiner Korpusnummer zitiert. Die Statusmarker [K], [B], [S] und
[SETZUNG] kennzeichnen den erreichten Ableitungsstand; spätere Korrekturen
und geschlossene Brücken sind im Korrekturregister Dok.~190 verzeichnet,
das laufend im Repository weitergeführt wird (siehe Band~5, Einleitung). Wo ein späteres Dokument eine Aussage dieses
Bandes präzisiert oder zurücknimmt, gilt die spätere Fassung. Die Dokumente
308 bis 310 sind ursprünglich als mehrkapitelige Abhandlungen gesetzt; im
Buch erscheinen ihre Kapitel als Abschnitte des jeweiligen Dokuments.

\subsection*{Stand}

Die Dokumente dieses Bandes tragen den Stand ihrer jeweiligen Fassung im
Korpus zum Oktober 2026.


	% --- part ---
	\part{Teil I --- Holografie, kosmologische Entartung und Formalismus-Korrespondenzen}
	\FFGFTdoc{Dok.~263 --- Fraktale Holografie: FFGFT und das holografische Prinzip}{../ch/263_Fraktale_Holografie_De_ch}
	\FFGFTdoc{Dok.~265 --- Formalismus-Korrespondenz: FFGFT und Blume-Maleev}{../ch/265_Formalismus_Korrespondenz_De_ch}
	\FFGFTdoc{Dok.~267 --- Kosmologische Entartung: $\Lambda$CDM und FFGFT}{../ch/267_Kosmologische_Entartung_De_ch}
	\FFGFTdoc{Dok.~268 --- CMB-Peaks aus dem T$^4$: Herleitung der Winkelspektrum-Verh\"altnisse}{../ch/268_CMB_Peaks_T4_De_ch}
	\FFGFTdoc{Dok.~269 --- FFGFT und RA/PMT: ein bidirektionaler Strukturvergleich}{../ch/269_FFGFT_RAPMT_Vergleich_De_ch}
	% --- part ---
	\part{Teil II --- Dimensionswechsel, HLV-Brücke und Skalenrekursion}
	\FFGFTdoc{Dok.~270 --- Dimensionswechsel in der FFGFT}{../ch/270_Projektionskette_T4_T0_De_ch}
	\FFGFTdoc{Dok.~271 --- FFGFT (T0) und Helix--Light--Vortex (HLV) im Vergleich}{../ch/271_FFGFT_HLV_Vergleich_De_ch}
	\FFGFTdoc{Dok.~272 --- FFGFT (T0) und HLV, eingeordnet nach Dot Theory}{../ch/272_FFGFT_HLV_DotTheory_Einordnung_De_ch}
	\FFGFTdoc{Dok.~274 --- FFGFT (T0) im IPI-Feld}{../ch/274_FFGFT_IPI_Feld_De_ch}
	\FFGFTdoc{Dok.~275 --- Von $\xi$ zu $\varphi$: Rekursive Skalenhierarchie}{../ch/275_xi_Rekursion_phi_De_ch}
	\FFGFTdoc{Dok.~276 --- Homologie-Test der HLV-Vortex-Hypothese}{../ch/276_HLV_Homologie_De_ch}
	\FFGFTdoc{Dok.~277 --- Wellenausbreitung und Energie}{../ch/277_Wellenausbreitung_Energie_De_ch}
	\FFGFTdoc{Dok.~279 --- Casimir, CMB und $H_0$ --- die gemeinsame $\xi$-Skala}{../ch/279_T0_Casimir_CMB_H0_De_ch}
	\FFGFTdoc{Dok.~280 --- Rotverschiebung in FFGFT: statisch, achromatisch}{../ch/280_T0_Rotverschiebung_Achromatisch_De_ch}
	\FFGFTdoc{Dok.~281 --- Der Goldene Schnitt in FFGFT}{../ch/281_Goldener_Schnitt_De_ch}
	\FFGFTdoc{Dok.~282 --- FFGFT im Hilbertraum --- Vom gemeinsamen Loop-Objekt zum Massenoperator}{../ch/282_FFGFT_HLV_CP_Teilbarkeit_De_ch}
	\FFGFTdoc{Dok.~283 --- FFGFT $\leftrightarrow$ HLV: die Brücke und die Gabelung}{../ch/283_FFGFT_HLV_Bruecke_De_ch}
	\FFGFTdoc{Dok.~284 --- Eine Formel-genaue Lesung des Spiral-Zeit-Neurodynamik-Modells}{../ch/284_FFGFT_Spiralzeit_Medin_De_ch}
	\FFGFTdoc{Dok.~285 --- Die einheitliche kompakte Sicht und die gerechnete 3-vs-6-Brücke}{../ch/285_FFGFT_HLV_Dimensionsbruecke_De_ch}
	% --- part ---
	\part{Teil III --- Dynamischer Sektor, Betrag und Phase, $\theta=2/9$}
	\FFGFTdoc{Dok.~286 --- Der dynamische Sektor: kinetische Energie als $\xi$-fraktaler Anteil des statischen}{../ch/286_Dynamischer_Sektor_Kinetik_De_ch}
	\FFGFTdoc{Dok.~287 --- Interne Darstellungs-Äquivalenzen}{../ch/287_Interne_Darstellungen_De_ch}
	\FFGFTdoc{Dok.~288 --- Nachrichtentechnische Rechenwege}{../ch/288_NT_Rechenwege_De_ch}
	\FFGFTdoc{Dok.~289 --- Magnitude/Phase-Karte der vier Fermion-Sektoren}{../ch/289_Magnitude_Phase_Sektoren_De_ch}
	\FFGFTdoc{Dok.~290 --- Folgen für die Informationsfrage: Betrag, Phase und der zweite Kanal}{../ch/290_Informationsfrage_Betrag_Phase_De_ch}
	\FFGFTdoc{Dok.~291 --- Der dynamische Ort von $\theta=2/9$: Betrag, Windung, Phase}{../ch/291_Dynamischer_Ort_Theta_2_9_De_ch}
	\FFGFTdoc{Dok.~292 --- Leptonen-Empirie-Check: die FFGFT-Relationen von den Daten aus}{../ch/292_Leptonen_Empirie_Check_De_ch}
	\FFGFTdoc{Dok.~293 --- Die ikosaedrische Herkunft von $\theta=2/9$: der zweite Zeuge}{../ch/293_Ikosaeder_Theta_2_9_De_ch}
	\FFGFTdoc{Dok.~294 --- FFGFT als Prüfinstanz für HLV: das geteilte $\varphi$-ikosaedrische Objekt}{../ch/294_HLV_Pruefrichtung_De_ch}
	% --- part ---
	\part{Teil IV --- Zeit und Information}
	\FFGFTdoc{Dok.~295 --- Die Zeit-Windung als logarithmische Spirale}{../ch/295_Zeit_Logspirale_De_ch}
	\FFGFTdoc{Dok.~296 --- Zeit, Offenheit und Invarianz}{../ch/296_Zeit_Offen_Invariant_De_ch}
	\FFGFTdoc{Dok.~297 --- HLV als mögliche Untermenge von FFGFT}{../ch/297_HLV_Untermenge_Kandidat_De_ch}
	\FFGFTdoc{Dok.~298 --- Das Standardmodell als dekompaktifizierte Projektion}{../ch/298_SM_Projektion_De_ch}
	\FFGFTdoc{Dok.~299 --- Die Einbettung $\iota$ — Arbeitsdokument und Forcing-Protokoll}{../ch/299_Einbettung_Iota_De_ch}
	\FFGFTdoc{Dok.~300 --- Spiral-Time als Auslesung und entprellter Comparator}{../ch/300_ZweiRollen_Kfrak_De_ch}
	\FFGFTdoc{Dok.~301 --- Der Umkehrtest --- Information als Ausgang, nicht als Grundeinheit}{../ch/301_Information_Umkehrtest_De_ch}
	\FFGFTdoc{Dok.~302 --- Elementarzellen, Ein-Bit-Kapazität und Partitionsinvarianz}{../ch/302_Elementarzellen_Spektrale_Antwort_De_ch}
	\FFGFTdoc{Dok.~303 --- Repräsentationale Identität über Rollenwechsel}{../ch/303_Repraesentationale_Identitaet_De_ch}
	\FFGFTdoc{Dok.~304 --- Der Gedächtnis-Reflex --- Warum erlebte Zeit das Gedächtnis voraussetzt}{../ch/304_Gedaechtnis_Reflex_De_ch}
	\FFGFTdoc{Dok.~305 --- Der Krümmungstreue-Diskriminator}{../ch/305_Kruemmungstreue_Diskriminator_De_ch}
	\FFGFTdoc{Dok.~306 --- Die native Zeit-Energie-Reziprozität}{../ch/306_Native_Zeit_Energie_Reziprozitaet_De_ch}
	\FFGFTdoc{Dok.~307 --- Zeit im Zustandsraum oder daneben --- Verortung der Zeit und Herkunft der Massen-Moden}{../ch/307_Zeit_im_Zustandsraum_De_ch}
	% --- part ---
	\part{Teil V --- Kosmischer Sektor und Resonanzgeometrie}
	\FFGFTdocDemote{Dok.~308 --- Der kosmische Sektor der FFGFT: Rotverschiebung, dunkler Sektor und geometrische Gravitation}{../ch/308_Lambda_Lesart_Artefakt_De_ch}
	\FFGFTdocDemote{Dok.~309 --- Das Skalenanker-Problem: $\Lambda$CDM und FFGFT im Vergleich}{../ch/309_Skalenanker_LCDM_FFGFT_De_ch}
	\FFGFTdocDemote{Dok.~310 --- Die Resonanzgeometrie --- als übergeordnete Referenz}{../ch/310_Resonanzgeometrie_Referenz_De_ch}
\end{document}
